\subsection{Exceptional partners and a shared failure-head budget}
\label{sec:partner-budget}

Let $D$ be a finite nonempty, strongly connected, all-deficient oriented
graph minimal under deletion of arbitrary nonempty arc sets. Fix $x$
with $P(x)\ne\varnothing$, and write
\[
 Q=N^{--}(x),\quad W=N^-(x),\quad C=\partial P(x),
 \quad X=\partial^2P(x),\quad H=Q\setminus C.
\]
Let $B=H\cap\mathcal M(x)$, $k=|B|$, and
$s=|H\cap N^+(x)|$, so $|H|=k+s$. Delete the $s$ arcs from $x$
into $H$, obtaining $\widehat D$; if $s=0$, it is the original graph.
Put
\[
 E=(\mathcal M(x)\setminus B)\,\dot\cup\,X,
 \qquad L=E\setminus N^{++}_{\widehat D}(x).
\]

\begin{lemma}[Shared head budget outside exceptional partners]
\label{lem:partner-head-budget}
Every $r\in B$ admits the convenient orientation $r\to x$.
Moreover
\[
 \bigcup_{r\in\mathcal M(x)\setminus B}
       \bigl(N^+(r)\cap U(x)\bigr)\subseteq L,
\]
and
\[
 |L|\le
 \begin{cases}
   \eta(x)-2k,&s=0,\\
   \eta(x)-g(x)-2k,&s>0.
 \end{cases}
\]
The bound concerns the union of distinct far heads, not the sum of
the sizes of the individual partner contributions.
\end{lemma}
\begin{proof}
The partition $V(D)=\{x\}\sqcup W\sqcup P(x)\sqcup C\sqcup H$
implies $\mathcal M(x)\setminus B\subseteq P(x)\cup C$.
All retained outneighbours of $x$ lie in $P(x)\cup C$.
Every arc leaving this union has its head in $X$, whereas a head
inside the union that is not an original outneighbour of $x$ is in
$\mathcal M(x)\setminus B$. Therefore the target envelope gives
$N^{++}_{\widehat D}(x)\subseteq E$, with the displayed union disjoint.
For $r\in\mathcal M(x)\setminus B$, any head of $r$ in $U(x)$ is
likewise in $E$. An original far head cannot become reachable after
deleting arcs, so it belongs to $L$.

The strict external-boundary inequality gives
$|X|\le |C|-1=|Q|-s-k-1$. Hence
\[
 |E|\le\mu(x)-k+|Q|-s-k-1
       =d_D^{++}(x)+\eta(x)-s-2k.
\]
If $s=0$, subtract $d_D^{++}(x)$. If $s>0$, the nonempty deletion
can create a Seymour vertex only at $x$, since all other first
neighbourhoods are unchanged and their second neighbourhoods cannot
grow. Minimality therefore gives
$d^{++}_{\widehat D}(x)\ge d_D^+(x)-s$; subtracting proves the second
bound. Finally every $r\in B$ lies outside both $P(x)$ and
$C=\partial P(x)$, so no member of $P(x)$ points to $r$.
All predecessors of $r$ thus reach $x$ within two steps, which is
exactly convenience of $r\to x$.
\end{proof}

For instance, if $\eta(x)=3$, $s=0$, and $k=1$, all nonexceptional
missing partners together contribute at most one far head. If
$\eta(x)=3$, $s>0$, and $k=1$, the bound forces $g(x)=1$ and $L$ empty;
every missing partner with an outneighbour in $U(x)$ must then be the
unique exceptional partner, directed conveniently into $x$.
Neither statement supplies a universal upper bound on $t(x)$.
